% 记得最后改名为 `AI 工具使用详情说明`
% 之所以这个文件用英文名是因为bibtex/biber似乎不支持含有中文的路径
\documentclass{MMCStyle}
\usepackage[multicol]{MyTemplate}

\input{Preamble for Pandoc.tex}

% \def\input@path{{dialog with AI}} % 可以设置多个路径

\NewDocumentEnvironment{dialog}{O{untitled} O{false}}{%
  \begin{solution}[#1]
  \def\multicolSwitch{#2}%
  \ifx\multicolSwitch\trueString
    \begin{multicols}{2}
  \fi
}{%
  \ifx\multicolSwitch\trueString
    \end{multicols}
  \fi
  \end{solution}
}

% 定义字符串比较辅助命令（放在导言区）
\def\trueString{true}
\def\falseString{false}

\bianhao{MC1111111}
\tihao{X}
\timu{参赛队 AI 工具使用详情说明}
\keyword{AI 工具使用详情}
% \bibliographystyle{gbt7714-numerical}
\addbibresource{ref.bib}


\begin{document}

\begin{abstract}
本文是参赛队（编号 MC1111111）的 AI 工具使用详情说明。文档将列出论文撰写过程中的 AI 工具的使用详情。
% 以及我们对 AI 生成内容进行处理后应用到论文的核心步骤。
\end{abstract}


% 取消目录页码
\tableofcontents\newpage
\clearpage
\pagenumbering{arabic}
\pagestyle{fancy}

\section{格式说明}

在后文中，深灰色背景的文字是我们提供给 AI 的提示词，框起来的具有浅灰色背景的文字是 AI 给的回答（均为 Markdown 语言，已经过渲染）。

给 AI （llm）提示词的时候注意不要让它使用 Unicode ，而是应该让它输出 Markdown 。因为你的 \LaTeX{} 编译器可能默认不支持某些 4 字节 UTF-8 字符。

注意 multicols 环境中不能使用浮动体，如果 pandoc 转换的表格用了 longtable 之类的东西，自己改或者让 AI 帮你改。参阅 \inline{雅可比矩阵介绍.tex} 和 \inline{雅可比矩阵.md} ，并留意 \inline{雅可比矩阵介绍.tex} 中被注释的表格。此外， multicols 环境内不支持 breakable 的盒子，这一点要么不要用 multicols （dialog 环境的可选参数为 true 表示使用 multicols 环境），要么不用我提供的 dialog 环境。

想改 markdown 的显示样式，自己改 \inline{Preamble for Pandoc.tex} 。有一些东西，比如 Markdown 的分割线 \inline{---} 是直接定义在 Pandoc 输出的 tex 里的，即
\begin{verbatim}
\begin{center}\rule{0.5\linewidth}{0.5pt}\end{center}
\end{verbatim}
想改的话可以改 \inline{ConvertMdToTex.py} 。

% 关于 \verb|\additionalSectoin| 系列命令产生的标题大小不受字号命令影响的问题，交给读者解决。这里提供两种思路：
% \begin{enumerate}
%   \item 重新定义 \lstinline|\additionalSection| 系列命令，使得它们不再依赖 \lstinline|\section*| 系列命令，而是简单地使用外部提供的字号并给必选参数加个 \lstinline|\textbf| 。
%   \item 
% \end{enumerate}

如果对这个 AI 工具使用详情有强迫症的话自己改 ConvertToTex.py ，反正我不想花那么多精力调这个东西。不值得。

\section{雅可比矩阵的介绍}

我们借助 Deepseek \cite{deepseek} 模型整理了关于雅可比矩阵的内容。

\begin{dialog}[介绍一下雅可比矩阵][false]
以下是关于\textbf{雅可比矩阵}的详细介绍，使用 Markdown 格式。

\begin{center}\rule{0.5\linewidth}{0.5pt}\end{center}

\additionalSection{雅可比矩阵（Jacobian
Matrix）}\label{ux96c5ux53efux6bd4ux77e9ux9635jacobian-matrix}

雅可比矩阵是多元微积分中的核心概念，可以看作\textbf{多元函数的一阶导数矩阵}。它描述了向量值函数在某一点附近如何随输入变化而变化。

\begin{center}\rule{0.5\linewidth}{0.5pt}\end{center}

\additionalSubsection{1. 定义}\label{ux5b9aux4e49}

设有一个函数

\[
\mathbf{f}: \mathbb{R}^n \to \mathbb{R}^m, \quad
\mathbf{y} = \mathbf{f}(\mathbf{x}), \quad
\mathbf{x} = (x_1, x_2, \dots, x_n),\quad
\mathbf{y} = (y_1, y_2, \dots, y_m)
\]

其中每个分量 \(y_i = f_i(\mathbf{x})\)。则 \textbf{雅可比矩阵} \(J\)
是一个 \(m \times n\) 的矩阵，其第 \(i\) 行第 \(j\) 列元素为：

\[
J_{ij} = \frac{\partial y_i}{\partial x_j}
\]

完整写出来：

\[
J = \begin{bmatrix}
\frac{\partial y_1}{\partial x_1} & \frac{\partial y_1}{\partial x_2} & \cdots & \frac{\partial y_1}{\partial x_n} \\
\frac{\partial y_2}{\partial x_1} & \frac{\partial y_2}{\partial x_2} & \cdots & \frac{\partial y_2}{\partial x_n} \\
\vdots & \vdots & \ddots & \vdots \\
\frac{\partial y_m}{\partial x_1} & \frac{\partial y_m}{\partial x_2} & \cdots & \frac{\partial y_m}{\partial x_n}
\end{bmatrix}
\]

\begin{center}\rule{0.5\linewidth}{0.5pt}\end{center}

\additionalSubsection{2.
几何与物理意义}\label{ux51e0ux4f55ux4e0eux7269ux7406ux610fux4e49}

\additionalSubsubsection{局部线性近似}\label{ux5c40ux90e8ux7ebfux6027ux8fd1ux4f3c}

在点 \(\mathbf{p}\) 附近，函数可以近似为：

\[
\mathbf{f}(\mathbf{p} + \Delta \mathbf{x}) \approx \mathbf{f}(\mathbf{p}) + J(\mathbf{p}) \cdot \Delta \mathbf{x}
\]

即雅可比矩阵给出了该点处的\textbf{最佳线性逼近}。

\additionalSubsubsection{\texorpdfstring{面积/体积伸缩（当
\(m = n\)）}{面积/体积伸缩（当 m = n）}}\label{ux9762ux79efux4f53ux79efux4f38ux7f29ux5f53-m-n}

如果 \(m = n\)，雅可比矩阵是方阵，其行列式称为 \textbf{雅可比行列式}
\(\det(J)\)。在坐标变换中，它表示局部体积（或面积）的缩放因子。例如极坐标变换：

\[
(x, y) = (r\cos\theta,\; r\sin\theta)
\] \[
\det(J) = r \quad \Rightarrow \quad dx\,dy = r\,dr\,d\theta
\]

\begin{center}\rule{0.5\linewidth}{0.5pt}\end{center}

\additionalSubsection{3. 特殊情形}\label{ux7279ux6b8aux60c5ux5f62}

% {\def\LTcaptype{none} % do not increment counter
% \begin{longtable}[]{@{}
%   >{\raggedright\arraybackslash}p{(\linewidth - 2\tabcolsep) * \real{0.2727}}
%   >{\raggedright\arraybackslash}p{(\linewidth - 2\tabcolsep) * \real{0.7273}}@{}}
% \toprule\noalign{}
% \begin{minipage}[b]{\linewidth}\raggedright
% 情形
% \end{minipage} & \begin{minipage}[b]{\linewidth}\raggedright
% 雅可比矩阵退化为
% \end{minipage} \\
% \midrule\noalign{}
% \endhead
% \bottomrule\noalign{}
% \endlastfoot
% \textbf{标量函数} \(f: \mathbb{R}^n \to \mathbb{R}\) & 行向量，即
% \textbf{梯度} \(\nabla f\) \\
% \textbf{向量值函数} \(\mathbf{f}: \mathbb{R} \to \mathbb{R}^m\) &
% 列向量，即各分量的普通导数 \\
% \textbf{规范映射} \(\mathbf{f}: \mathbb{R}^n \to \mathbb{R}^n\)
% 且为线性变换 & 常数矩阵（即变换矩阵本身） \\
% \end{longtable}
% }

\begin{tabular}{|l|p{8cm}|}
\hline
\textbf{领域} & \textbf{用途} \\
\hline
优化 & 梯度下降法（雅可比→梯度） \\
\hline
机器人学 & 关节速度与末端执行器速度之间的映射（运动学雅可比） \\
\hline
几何/物理 & 坐标变换中的体积元（例：极坐标 $\diff x\diff y = r \diff r \diff\theta$，雅可比行列式为 $r$） \\
\hline
数值计算 & 牛顿法解方程组（雅可比矩阵用于线性化） \\
\hline
深度学习 & 反向传播本质上就是在计算雅可比矩阵的乘积 \\
\hline
\end{tabular}

\begin{center}\rule{0.5\linewidth}{0.5pt}\end{center}

\additionalSubsection{4. 典型应用}\label{ux5178ux578bux5e94ux7528}

% {\def\LTcaptype{none} % do not increment counter
% \begin{longtable}[]{@{}ll@{}}
% \toprule\noalign{}
% 领域 & 用途 \\
% \midrule\noalign{}
% \endhead
% \bottomrule\noalign{}
% \endlastfoot
% \textbf{优化} & 梯度下降法中，梯度即雅可比（标量情形） \\
% \textbf{机器人学} & 运动学雅可比：关节速度 → 末端执行器速度 \\
% \textbf{几何与物理} & 坐标变换中的体积元（雅可比行列式） \\
% \textbf{数值计算} & 牛顿法解非线性方程组：利用雅可比线性化 \\
% \textbf{深度学习} & 反向传播算法本质上是雅可比矩阵的链式乘积 \\
% \end{longtable}
% }

\begin{tabular}{ll}
\toprule
\textbf{领域} & \textbf{用途} \\
\midrule
优化 & 梯度下降法中，梯度即雅可比（标量情形） \\
机器人学 & 运动学雅可比：关节速度 $rightarrow$ 末端执行器速度 \\
几何与物理 & 坐标变换中的体积元（雅可比行列式） \\
数值计算 & 牛顿法解非线性方程组：利用雅可比线性化 \\
深度学习 & 反向传播算法本质上是雅可比矩阵的链式乘积 \\
\bottomrule
\end{tabular}

\begin{center}\rule{0.5\linewidth}{0.5pt}\end{center}

\additionalSubsection{5. 例子}\label{ux4f8bux5b50}

\additionalSubsubsection{\texorpdfstring{例1：极坐标变换（\(\mathbb{R}^2 \to \mathbb{R}^2\)）}{例1：极坐标变换（\textbackslash mathbb\{R\}\^{}2 \textbackslash to \textbackslash mathbb\{R\}\^{}2）}}\label{ux4f8b1ux6781ux5750ux6807ux53d8ux6362mathbbr2-to-mathbbr2}

\[
f(r,\theta) = (r\cos\theta,\; r\sin\theta)
\] \[
J = \begin{bmatrix}
\cos\theta & -r\sin\theta \\
\sin\theta &  r\cos\theta
\end{bmatrix},\quad \det(J) = r
\]

\additionalSubsubsection{\texorpdfstring{例2：简单非线性映射（\(\mathbb{R}^2 \to \mathbb{R}^2\)）}{例2：简单非线性映射（\textbackslash mathbb\{R\}\^{}2 \textbackslash to \textbackslash mathbb\{R\}\^{}2）}}\label{ux4f8b2ux7b80ux5355ux975eux7ebfux6027ux6620ux5c04mathbbr2-to-mathbbr2}

\[
f(x,y) = (x^2 y,\; x + y^3)
\] \[
J = \begin{bmatrix}
2xy & x^2 \\
1   & 3y^2
\end{bmatrix}
\]

\begin{center}\rule{0.5\linewidth}{0.5pt}\end{center}

\additionalSubsection{6.
与黑塞矩阵的区别}\label{ux4e0eux9ed1ux585eux77e9ux9635ux7684ux533aux522b}

% {\def\LTcaptype{none} % do not increment counter
% \begin{longtable}[]{@{}lll@{}}
% \toprule\noalign{}
% 矩阵 & 阶数 & 含义 \\
% \midrule\noalign{}
% \endhead
% \bottomrule\noalign{}
% \endlastfoot
% \textbf{雅可比矩阵} & 一阶偏导 & 局部线性变化率 \\
% \textbf{黑塞矩阵} & 二阶偏导 & 局部曲率（仅用于标量函数） \\
% \end{longtable}
% }

\begin{tabular}{lll}
\toprule
矩阵 & 阶数 & 含义 \\
\midrule
雅可比矩阵 & 一阶偏导 & 局部线性变化率 \\
黑塞矩阵 & 二阶偏导 & 局部曲率（仅用于标量函数） \\
\bottomrule
\end{tabular}

\begin{center}\rule{0.5\linewidth}{0.5pt}\end{center}

\additionalSubsection{总结}\label{ux603bux7ed3}

\begin{itemize}
\tightlist
\item
  雅可比矩阵是多元向量值函数的\textbf{一阶导数矩阵}
\item
  提供局部线性近似，是微分、几何、数值计算与工程中的基础工具
\item
  当输出为标量时退化为梯度；当输入输出维度相同时，其行列式表示空间伸缩比
\end{itemize}

如果你对某个具体应用（如机器人运动学或坐标变换）有更深入的问题，欢迎继续提问！

\end{dialog}

\section{雅可比矩阵的介绍2}

我们借助 Deepseek \cite{deepseek} 模型整理了关于雅可比矩阵的内容。

\begin{dialog}[介绍一下雅可比矩阵][true]
\zihao{7}
以下是关于\textbf{雅可比矩阵}的详细介绍，使用 Markdown 格式。

\begin{center}\rule{0.5\linewidth}{0.5pt}\end{center}

\additionalSection{雅可比矩阵（Jacobian
Matrix）}\label{ux96c5ux53efux6bd4ux77e9ux9635jacobian-matrix}

雅可比矩阵是多元微积分中的核心概念，可以看作\textbf{多元函数的一阶导数矩阵}。它描述了向量值函数在某一点附近如何随输入变化而变化。

\begin{center}\rule{0.5\linewidth}{0.5pt}\end{center}

\additionalSubsection{1. 定义}\label{ux5b9aux4e49}

设有一个函数

\[
\mathbf{f}: \mathbb{R}^n \to \mathbb{R}^m, \quad
\mathbf{y} = \mathbf{f}(\mathbf{x}), \quad
\mathbf{x} = (x_1, x_2, \dots, x_n),\quad
\mathbf{y} = (y_1, y_2, \dots, y_m)
\]

其中每个分量 \(y_i = f_i(\mathbf{x})\)。则 \textbf{雅可比矩阵} \(J\)
是一个 \(m \times n\) 的矩阵，其第 \(i\) 行第 \(j\) 列元素为：

\[
J_{ij} = \frac{\partial y_i}{\partial x_j}
\]

完整写出来：

\[
J = \begin{bmatrix}
\frac{\partial y_1}{\partial x_1} & \frac{\partial y_1}{\partial x_2} & \cdots & \frac{\partial y_1}{\partial x_n} \\
\frac{\partial y_2}{\partial x_1} & \frac{\partial y_2}{\partial x_2} & \cdots & \frac{\partial y_2}{\partial x_n} \\
\vdots & \vdots & \ddots & \vdots \\
\frac{\partial y_m}{\partial x_1} & \frac{\partial y_m}{\partial x_2} & \cdots & \frac{\partial y_m}{\partial x_n}
\end{bmatrix}
\]

\begin{center}\rule{0.5\linewidth}{0.5pt}\end{center}

\additionalSubsection{2.
几何与物理意义}\label{ux51e0ux4f55ux4e0eux7269ux7406ux610fux4e49}

\additionalSubsubsection{局部线性近似}\label{ux5c40ux90e8ux7ebfux6027ux8fd1ux4f3c}

在点 \(\mathbf{p}\) 附近，函数可以近似为：

\[
\mathbf{f}(\mathbf{p} + \Delta \mathbf{x}) \approx \mathbf{f}(\mathbf{p}) + J(\mathbf{p}) \cdot \Delta \mathbf{x}
\]

即雅可比矩阵给出了该点处的\textbf{最佳线性逼近}。

\additionalSubsubsection{\texorpdfstring{面积/体积伸缩（当
\(m = n\)）}{面积/体积伸缩（当 m = n）}}\label{ux9762ux79efux4f53ux79efux4f38ux7f29ux5f53-m-n}

如果 \(m = n\)，雅可比矩阵是方阵，其行列式称为 \textbf{雅可比行列式}
\(\det(J)\)。在坐标变换中，它表示局部体积（或面积）的缩放因子。例如极坐标变换：

\[
(x, y) = (r\cos\theta,\; r\sin\theta)
\] \[
\det(J) = r \quad \Rightarrow \quad dx\,dy = r\,dr\,d\theta
\]

\begin{center}\rule{0.5\linewidth}{0.5pt}\end{center}

\additionalSubsection{3. 特殊情形}\label{ux7279ux6b8aux60c5ux5f62}

% {\def\LTcaptype{none} % do not increment counter
% \begin{longtable}[]{@{}
%   >{\raggedright\arraybackslash}p{(\linewidth - 2\tabcolsep) * \real{0.2727}}
%   >{\raggedright\arraybackslash}p{(\linewidth - 2\tabcolsep) * \real{0.7273}}@{}}
% \toprule\noalign{}
% \begin{minipage}[b]{\linewidth}\raggedright
% 情形
% \end{minipage} & \begin{minipage}[b]{\linewidth}\raggedright
% 雅可比矩阵退化为
% \end{minipage} \\
% \midrule\noalign{}
% \endhead
% \bottomrule\noalign{}
% \endlastfoot
% \textbf{标量函数} \(f: \mathbb{R}^n \to \mathbb{R}\) & 行向量，即
% \textbf{梯度} \(\nabla f\) \\
% \textbf{向量值函数} \(\mathbf{f}: \mathbb{R} \to \mathbb{R}^m\) &
% 列向量，即各分量的普通导数 \\
% \textbf{规范映射} \(\mathbf{f}: \mathbb{R}^n \to \mathbb{R}^n\)
% 且为线性变换 & 常数矩阵（即变换矩阵本身） \\
% \end{longtable}
% }

\begin{tabular}{|l|p{8cm}|}
\hline
\textbf{领域} & \textbf{用途} \\
\hline
优化 & 梯度下降法（雅可比→梯度） \\
\hline
机器人学 & 关节速度与末端执行器速度之间的映射（运动学雅可比） \\
\hline
几何/物理 & 坐标变换中的体积元（例：极坐标 $\diff x\diff y = r \diff r \diff\theta$，雅可比行列式为 $r$） \\
\hline
数值计算 & 牛顿法解方程组（雅可比矩阵用于线性化） \\
\hline
深度学习 & 反向传播本质上就是在计算雅可比矩阵的乘积 \\
\hline
\end{tabular}

\begin{center}\rule{0.5\linewidth}{0.5pt}\end{center}

\additionalSubsection{4. 典型应用}\label{ux5178ux578bux5e94ux7528}

% {\def\LTcaptype{none} % do not increment counter
% \begin{longtable}[]{@{}ll@{}}
% \toprule\noalign{}
% 领域 & 用途 \\
% \midrule\noalign{}
% \endhead
% \bottomrule\noalign{}
% \endlastfoot
% \textbf{优化} & 梯度下降法中，梯度即雅可比（标量情形） \\
% \textbf{机器人学} & 运动学雅可比：关节速度 → 末端执行器速度 \\
% \textbf{几何与物理} & 坐标变换中的体积元（雅可比行列式） \\
% \textbf{数值计算} & 牛顿法解非线性方程组：利用雅可比线性化 \\
% \textbf{深度学习} & 反向传播算法本质上是雅可比矩阵的链式乘积 \\
% \end{longtable}
% }

\begin{tabular}{ll}
\toprule
\textbf{领域} & \textbf{用途} \\
\midrule
优化 & 梯度下降法中，梯度即雅可比（标量情形） \\
机器人学 & 运动学雅可比：关节速度 $rightarrow$ 末端执行器速度 \\
几何与物理 & 坐标变换中的体积元（雅可比行列式） \\
数值计算 & 牛顿法解非线性方程组：利用雅可比线性化 \\
深度学习 & 反向传播算法本质上是雅可比矩阵的链式乘积 \\
\bottomrule
\end{tabular}

\begin{center}\rule{0.5\linewidth}{0.5pt}\end{center}

\additionalSubsection{5. 例子}\label{ux4f8bux5b50}

\additionalSubsubsection{\texorpdfstring{例1：极坐标变换（\(\mathbb{R}^2 \to \mathbb{R}^2\)）}{例1：极坐标变换（\textbackslash mathbb\{R\}\^{}2 \textbackslash to \textbackslash mathbb\{R\}\^{}2）}}\label{ux4f8b1ux6781ux5750ux6807ux53d8ux6362mathbbr2-to-mathbbr2}

\[
f(r,\theta) = (r\cos\theta,\; r\sin\theta)
\] \[
J = \begin{bmatrix}
\cos\theta & -r\sin\theta \\
\sin\theta &  r\cos\theta
\end{bmatrix},\quad \det(J) = r
\]

\additionalSubsubsection{\texorpdfstring{例2：简单非线性映射（\(\mathbb{R}^2 \to \mathbb{R}^2\)）}{例2：简单非线性映射（\textbackslash mathbb\{R\}\^{}2 \textbackslash to \textbackslash mathbb\{R\}\^{}2）}}\label{ux4f8b2ux7b80ux5355ux975eux7ebfux6027ux6620ux5c04mathbbr2-to-mathbbr2}

\[
f(x,y) = (x^2 y,\; x + y^3)
\] \[
J = \begin{bmatrix}
2xy & x^2 \\
1   & 3y^2
\end{bmatrix}
\]

\begin{center}\rule{0.5\linewidth}{0.5pt}\end{center}

\additionalSubsection{6.
与黑塞矩阵的区别}\label{ux4e0eux9ed1ux585eux77e9ux9635ux7684ux533aux522b}

% {\def\LTcaptype{none} % do not increment counter
% \begin{longtable}[]{@{}lll@{}}
% \toprule\noalign{}
% 矩阵 & 阶数 & 含义 \\
% \midrule\noalign{}
% \endhead
% \bottomrule\noalign{}
% \endlastfoot
% \textbf{雅可比矩阵} & 一阶偏导 & 局部线性变化率 \\
% \textbf{黑塞矩阵} & 二阶偏导 & 局部曲率（仅用于标量函数） \\
% \end{longtable}
% }

\begin{tabular}{lll}
\toprule
矩阵 & 阶数 & 含义 \\
\midrule
雅可比矩阵 & 一阶偏导 & 局部线性变化率 \\
黑塞矩阵 & 二阶偏导 & 局部曲率（仅用于标量函数） \\
\bottomrule
\end{tabular}

\begin{center}\rule{0.5\linewidth}{0.5pt}\end{center}

\additionalSubsection{总结}\label{ux603bux7ed3}

\begin{itemize}
\tightlist
\item
  雅可比矩阵是多元向量值函数的\textbf{一阶导数矩阵}
\item
  提供局部线性近似，是微分、几何、数值计算与工程中的基础工具
\item
  当输出为标量时退化为梯度；当输入输出维度相同时，其行列式表示空间伸缩比
\end{itemize}

如果你对某个具体应用（如机器人运动学或坐标变换）有更深入的问题，欢迎继续提问！

\end{dialog}

\clearpage
\phantomsection
\addcontentsline{toc}{section}{参考文献}
% \bibliography{ref}
\printbibliography[title={参考文献}] % 输出列表

\end{document}

